%====================================================================%
% Supplementary tables                                               %
% Generated by src/make_tables.py from results/. Do not edit by hand. %
% To change a number, re-run the analysis and then this script.      %
% Journal rules: caption above (<=15 words), legend below, no colour, %
% no thousands separators.                                           %
%====================================================================%

\begin{table}[h]
\caption{Sensitivity of the primary result to three data-curation decisions}\label{tab:curation}
\begin{tabular}{@{}llccccc@{}}
\toprule
Dataset & Condition & $n$ & ECE & Control ECE & Ratio & Reliability share \\
\midrule
DILIrank & Primary analysis & 796 & 0.4384 & 0.1300 & 3.37 & 86.1\% \\
 & Corrected approval years & 796 & 0.4386 & 0.1300 & 3.37 & 86.2\% \\
 & Corrected years, one exclusion & 795 & 0.4426 & 0.1201 & 3.68 & 84.7\% \\
 & ATC class V removed & 752 & 0.4574 & 0.1329 & 3.44 & 91.2\% \\
 & Stereoisomers merged (primary key) & 796 & 0.4384 & 0.2500 & 1.75 & 83.1\% \\
 & Stereoisomers separated & 797 & 0.4325 & 0.2462 & 1.76 & 84.6\% \\
\midrule
DICTrank & Primary analysis & 919 & 0.2702 & 0.1377 & 1.96 & 82.5\% \\
 & Corrected approval years & 919 & 0.2816 & 0.1377 & 2.05 & 80.3\% \\
 & Corrected years, one exclusion & 919 & 0.2816 & 0.1377 & 2.05 & 80.3\% \\
 & ATC class V removed & 875 & 0.2547 & 0.1375 & 1.85 & 86.7\% \\
 & Stereoisomers merged (primary key) & 919 & 0.2702 & 0.1831 & 1.48 & 62.3\% \\
 & Stereoisomers separated & 930 & 0.2724 & 0.1922 & 1.42 & 77.3\% \\
\botrule
\end{tabular}
\footnotetext{\textbf{The control ECE column is the denominator of the ratio in the same row and differs between rows}: the approval-year and ATC rows use a size-matched random control re-drawn for each analysis set, while the stereoisomer rows use the prevalence-matched control of Table~\ref{tab:prev}, so ratios are comparable within a block and not across blocks. The corrected-years row applies the seven approval-year corrections identified by the audit, one of which the manual review judged insufficiently supported; the next row additionally excludes phenobarbital sodium. Reliability share is the Murphy decomposition share of Table~\ref{tab:brier}. Across all 12 conditions the ratio stays above 1 (range 1.42--3.68) and the reliability share above half (range 62.3\%--91.2\%); the lowest share, 62.3\%, is the DICTrank primary key, which is why no claim of a dominant calibration component is made for that dataset.}
\end{table}

\begin{table}[h]
\caption{Structure alone separates pre- and post-2010 approvals}\label{tab:era}
\begin{tabular}{@{}lccccccc@{}}
\toprule
Dataset & $n$ (post-2010) & Era AUC & Null mean (SD) & Null 2.5--97.5\% & At or above & $p$ & $z$ \\
\midrule
DILIrank & 796 (156) & 0.6835 & 0.5003 (0.0359) & 0.4297--0.5707 & 0/2000 & 0.0005 & 5.11 \\
DICTrank & 919 (267) & 0.7917 & 0.4982 (0.0293) & 0.4394--0.5528 & 0/2000 & 0.0005 & 10.02 \\
\botrule
\end{tabular}
\footnotetext{A classifier trained on the same representation and learner to predict whether a compound was approved from 2010 onwards, scored by five-fold cross-validated AUC. The null distribution randomises the era label 2000 times. An AUC of 0.5 would mean the two halves are structurally indistinguishable. No permutation reached the observed value in either dataset, so $p$ sits at the 0.0005 bound and the true value is smaller. \textbf{The $z$ column is not a test}: the null distribution is not known to be normal, so it must not be converted into a $p$ value; it is given only because the bounded $p$ cannot show how far the observation lies from the null. Applying the identical procedure to random split labels of the same sizes gives 0.4991 and 0.4924, that is chance, which rules out the objection that the classifier would separate any partition. The six RDKit descriptors alone, without fingerprints, still separate the eras (0.5951 and 0.6805), so the shift is not an artefact of fingerprint sparsity (Table~\ref{tab:desc}).}
\end{table}

\begin{table}[h]
\caption{Median RDKit descriptors in the training and held-out halves}\label{tab:desc}
\begin{tabular}{@{}lcccc@{}}
\toprule
Descriptor & DILIrank train & DILIrank test & DICTrank train & DICTrank test \\
\midrule
Molecular weight & 336.2 & 441.0 (+104.8) & 340.0 & 450.3 (+110.4) \\
Calculated logP & 2.26 & 2.99 (+0.73) & 2.19 & 3.43 (+1.24) \\
Topological polar surface area & 75.4 & 98.5 (+23.0) & 76.1 & 96.9 (+20.8) \\
H-bond acceptors & 4 & 6 (+2) & 4 & 6 (+2) \\
Rotatable bonds & 4 & 6 (+2) & 5 & 6 (+1) \\
H-bond donors & 2 & 2 (+0) & 2 & 2 (+0) \\
\botrule
\end{tabular}
\footnotetext{Medians of the six descriptors used by the primary model, with the signed change from training to held-out in parentheses. Compounds approved from 2010 onwards are larger, more lipophilic and have greater polar surface area in both datasets. The distribution shift the temporal split creates is therefore real at the level of structure, which is what Table~\ref{tab:era} measures as a single number. \textbf{Existence of shift is not evidence that its magnitude drives the calibration collapse}; Table~\ref{tab:nontemporal} tests that and does not support it.}
\end{table}

\begin{table}[h]
\caption{Dependence between the two datasets on their shared held-out compounds}\label{tab:overlap}
\begin{tabular}{@{}lcccc@{}}
\toprule
Quantity & $n$ & Pearson $r$ & 95\% CI & Spearman \\
\midrule
Signed residual $y-p$ & 135 & +0.2340 & [0.075--0.384] & +0.1498 \\
Absolute residual $|y-p|$ & 135 & +0.0888 & [-0.079, 0.253] & +0.0483 \\
Label $y$ (phi coefficient) & 135 & +0.1600 & [-0.011, 0.324] & +0.1600 \\
Predicted probability $p$ & 135 & +0.0981 & [-0.127, 0.294] & -0.0588 \\
\botrule
\end{tabular}
\footnotetext{The two datasets share compounds, so a reviewer may ask whether they constitute independent replication. Of the held-out sets, 135 compounds appear in both. The signed residual is the quantity that matters, because ECE is built from signed mean residuals within bins; its correlation is +0.2340 [0.075--0.384], which excludes zero. Averaging the two estimates therefore reduces variance to 0.62 of one estimate, equivalent to 1.62 independent datasets, or 1.44 if the upper confidence bound is used. The error magnitudes are a weaker signal, +0.0888 [-0.079, 0.253], and the labels themselves are nearly independent, phi = +0.1600 on these same compounds. \textbf{The two results are not fully independent replication and this table quantifies by how much}; Table~\ref{tab:nontemporal} reports the disjoint-subset replication separately. The equal-variance approximation behind the effective count is stated rather than tested.}
\end{table}

\begin{table}[h]
\caption{Refusing to predict on the least familiar compounds does not help}\label{tab:adfilter}
\begin{tabular}{@{}lccccccc@{}}
\toprule
Dataset & Fraction dropped & Coverage & $n$ kept & ECE & AUROC & ECE, random drop & Difference \\
\midrule
DILIrank & 0\% & 1.000 & 156 & 0.4384 & 0.6476 & 0.4384 & -0.0000 \\
 & 10\% & 0.897 & 140 & 0.4333 & 0.6556 & 0.4413 & -0.0079 \\
 & 20\% & 0.801 & 125 & 0.4288 & 0.6654 & 0.4434 & -0.0147 \\
 & 30\% & 0.699 & 109 & 0.4650 & 0.6531 & 0.4447 & +0.0204 \\
 & 40\% & 0.603 & 94 & 0.4741 & 0.6604 & 0.4469 & +0.0272 \\
 & 50\% & 0.500 & 78 & 0.4700 & 0.6557 & 0.4521 & +0.0179 \\
\midrule
DICTrank & 0\% & 1.000 & 267 & 0.2702 & 0.6198 & 0.2702 & +0.0000 \\
 & 10\% & 0.899 & 240 & 0.2965 & 0.6160 & 0.2726 & +0.0239 \\
 & 20\% & 0.801 & 214 & 0.2952 & 0.6426 & 0.2749 & +0.0204 \\
 & 30\% & 0.700 & 187 & 0.2856 & 0.6570 & 0.2766 & +0.0090 \\
 & 40\% & 0.599 & 160 & 0.2500 & 0.6846 & 0.2792 & -0.0292 \\
 & 50\% & 0.498 & 133 & 0.2785 & 0.7058 & 0.2820 & -0.0034 \\
\botrule
\end{tabular}
\footnotetext{Compounds are dropped in order of increasing nearest-neighbour similarity to the training set, which is the applicability-domain filter a reviewer would propose. The control drops the same number at random, and because the predictions are fixed neither arm refits a model. The difference column must be negative for the filter to be doing anything beyond shrinking the sample. \textbf{It is positive in 6 of the 10 non-zero settings}, and where negative it is small, so the sign is not stable. An applicability-domain filter is therefore not a remedy for the calibration loss reported in Table~\ref{tab:split}.}
\end{table}

\begin{table}[h]
\caption{The headline ratio holds across every model and representation}\label{tab:gridr3}
\begin{tabular}{@{}lllccccc@{}}
\toprule
Dataset & Representation & Model & Temporal ECE & Control ECE & Ratio & $\Delta$reliability & $\Delta$resolution \\
\midrule
DILIrank & ECFP4+desc & XGBoost & 0.4384 & 0.3194 & 1.37 & +0.0985 & -0.0228 \\
 & ECFP4+desc & RandomForest & 0.3584 & 0.3175 & 1.13 & +0.0228 & -0.0230 \\
 & ECFP4+desc & LogisticReg & 0.4265 & 0.3502 & 1.22 & +0.0681 & +0.0009 \\
 & ECFP4+desc & LightGBM & 0.5059 & 0.3592 & 1.41 & +0.1203 & -0.0180 \\
 & ECFP4 & XGBoost & 0.3977 & 0.3047 & 1.31 & +0.0763 & -0.0218 \\
 & ECFP4 & RandomForest & 0.3530 & 0.3182 & 1.11 & +0.0179 & -0.0254 \\
 & ECFP4 & LogisticReg & 0.4140 & 0.3501 & 1.18 & +0.0585 & -0.0087 \\
 & ECFP4 & LightGBM & 0.4641 & 0.3257 & 1.42 & +0.1131 & -0.0166 \\
 & MACCS & XGBoost & 0.4372 & 0.3229 & 1.35 & +0.0969 & -0.0256 \\
 & MACCS & RandomForest & 0.3704 & 0.2925 & 1.27 & +0.0548 & -0.0193 \\
 & MACCS & LogisticReg & 0.4016 & 0.3130 & 1.28 & +0.0683 & -0.0175 \\
 & MACCS & LightGBM & 0.4345 & 0.3283 & 1.32 & +0.0911 & -0.0326 \\
 & RDKit desc & XGBoost & 0.4557 & 0.3672 & 1.24 & +0.0689 & -0.0232 \\
 & RDKit desc & RandomForest & 0.3900 & 0.3168 & 1.23 & +0.0567 & -0.0099 \\
 & RDKit desc & LogisticReg & 0.3907 & 0.3304 & 1.18 & +0.0358 & -0.0014 \\
 & RDKit desc & LightGBM & 0.4644 & 0.3754 & 1.24 & +0.0724 & -0.0213 \\
\midrule
DICTrank & ECFP4+desc & XGBoost & 0.2702 & 0.2246 & 1.20 & +0.0232 & -0.0186 \\
 & ECFP4+desc & RandomForest & 0.1884 & 0.1592 & 1.18 & +0.0093 & -0.0163 \\
 & ECFP4+desc & LogisticReg & 0.3253 & 0.2807 & 1.16 & +0.0243 & -0.0151 \\
 & ECFP4+desc & LightGBM & 0.3363 & 0.2756 & 1.22 & +0.0346 & -0.0165 \\
 & ECFP4 & XGBoost & 0.2585 & 0.2101 & 1.23 & +0.0174 & -0.0119 \\
 & ECFP4 & RandomForest & 0.1848 & 0.1577 & 1.17 & +0.0096 & -0.0114 \\
 & ECFP4 & LogisticReg & 0.3192 & 0.2817 & 1.13 & +0.0222 & -0.0184 \\
 & ECFP4 & LightGBM & 0.3110 & 0.2573 & 1.21 & +0.0312 & -0.0088 \\
 & MACCS & XGBoost & 0.2701 & 0.2294 & 1.18 & +0.0275 & -0.0147 \\
 & MACCS & RandomForest & 0.1718 & 0.1381 & 1.24 & +0.0069 & -0.0242 \\
 & MACCS & LogisticReg & 0.2567 & 0.2249 & 1.14 & +0.0176 & -0.0059 \\
 & MACCS & LightGBM & 0.3191 & 0.2573 & 1.24 & +0.0330 & -0.0141 \\
 & RDKit desc & XGBoost & 0.2667 & 0.2457 & 1.09 & +0.0135 & -0.0116 \\
 & RDKit desc & RandomForest & 0.1877 & 0.1565 & 1.20 & +0.0096 & -0.0197 \\
 & RDKit desc & LogisticReg & 0.2103 & 0.1595 & 1.32 & +0.0152 & -0.0072 \\
 & RDKit desc & LightGBM & 0.3105 & 0.2768 & 1.12 & +0.0263 & -0.0092 \\
\botrule
\end{tabular}
\footnotetext{The control matches the prevalence of both the training and the validation set, which is the most conservative of the three designs used in this work. Each entry averages 15 control repetitions. Temporal ECE was verified to reproduce the grid reported elsewhere to within $10^{-6}$ in all 32 cells. \textbf{All 32 ratios exceed one}, spanning 1.09 to 1.42 with a median of 1.22, and the medians by model (1.18--1.24) and by representation (1.20--1.25) lie in a narrow band, so no family is immune. The share of the Brier increase attributable to reliability is not reported here: its denominator passes through zero at this number of repetitions and the percentage is therefore uninterpretable. The sign comparison is reported instead, and reliability worsens in 32 of 32 cells, exceeding the absolute change in resolution in 26.}
\end{table}

\begin{table}[h]
\caption{Calibration loss appears in every approval window after control}\label{tab:horizon}
\begin{tabular}{@{}llcccccc@{}}
\toprule
Dataset & Approval window & $n$ & Prevalence & ECE & Ratio (R1) & Ratio (R3) & 95\% CI \\
\midrule
DILIrank & 2010-2013 & 50 & 0.5200 & 0.3136 & 1.81 & 1.41 & [0.949--2.049] \\
 & 2014-2017 & 49 & 0.2857 & 0.4385 & 2.52 & 1.24 & [0.956--1.636] \\
 & 2018-2022 & 57 & 0.1404 & 0.6158 & 3.67 & 1.39 & [1.155--1.619] \\
\midrule
DICTrank & 2010-2013 & 65 & 0.7231 & 0.2078 & 1.32 & 1.20 & [0.692--1.830] \\
 & 2014-2017 & 82 & 0.5732 & 0.3162 & 2.18 & 1.20 & [0.881--1.649] \\
 & 2018-2025 & 120 & 0.5833 & 0.3074 & 2.27 & 1.24 & [0.968--1.632] \\
\botrule
\end{tabular}
\footnotetext{The validation set is split by approval window and each window is compared with controls of the same size and prevalence; the R3 control matches the prevalence of both the training and the validation side, so the only remaining difference from the temporal split is the axis. Intervals are 2.5--97.5 percentiles of a bootstrap of the window, with the predictions held fixed. \textbf{All 6 window ratios exceed one}, but only 1 interval excludes one, because windows hold 49 to 120 compounds. Absolute ECE rises across windows in both datasets and so does the uncontrolled R1 ratio, yet the change in the R3 ratio from the first window to the last is -0.0212 and +0.0348, smaller than the mean half-width of the intervals at the two ends (0.391 and 0.450). The rising absolute trend therefore tracks the falling prevalence rather than the prediction horizon, and a horizon effect can be neither claimed nor ruled out at this sample size.}
\end{table}

\begin{table}[h]
\caption{Nine quantification estimators do not recover the validation prevalence}\label{tab:prior2}
\begin{tabular}{@{}lllccccc@{}}
\toprule
Dataset & Estimator & Information used & $\hat{\pi}$ & Error & 95\% CI & ECE after & Recovery \\
\midrule
DILIrank & Q1 CC & none & 0.7821 & +0.4744 & [0.718--0.846] & 0.4966 & -12.5\% \\
 & Q2 ACC & training labels only & 0.7568 & +0.4491 & [0.591--0.906] & 0.4933 & -11.7\% \\
 & Q3 PACC & training labels only & 0.7277 & +0.4201 & [0.572--0.877] & 0.4705 & -6.0\% \\
 & Q4 SLD/EM, 전체 학습 & none & 0.8235 & +0.5158 & [0.731--0.907] & 0.5158 & -17.4\% \\
 & Q4b SLD/EM, 보정셋 제외 학습 & none & 0.8210 & +0.5133 & [0.729--0.900] & 0.5145 & -17.0\% \\
 & Q5 SLD/EM + Platt 보정 & training labels only & 0.8514 & +0.5437 & [0.725--0.968] & 0.5314 & -21.3\% \\
 & Q6 ACC/BBSE-hard & training labels only & 0.7568 & +0.4491 & [0.591--0.923] & 0.4933 & -11.7\% \\
 & Q7 BBSE-soft & training labels only & 0.7277 & +0.4201 & [0.568--0.865] & 0.4705 & -6.0\% \\
 & Q8 HDy & training labels only & 0.7300 & +0.4223 & [0.620--0.870] & 0.4712 & -6.2\% \\
 & Q9 MS & training labels only & 0.7298 & +0.4221 & [0.586--0.868] & 0.4711 & -6.1\% \\
 & Q0 참값 & validation labels & 0.3077 & +0.0000 & [0.237--0.385] & 0.2959 & 37.8\% \\
\midrule
DICTrank & Q1 CC & none & 0.9176 & +0.3034 & [0.884--0.951] & 0.3312 & -23.5\% \\
 & Q2 ACC & training labels only & 0.8945 & +0.2803 & [0.000--1.000] & 0.3179 & -18.2\% \\
 & Q3 PACC & training labels only & 0.5257 & -0.0886 & [0.119--0.876] & 0.2210 & 20.7\% \\
 & Q4 SLD/EM, 전체 학습 & none & 0.9550 & +0.3408 & [0.897--1.000] & 0.3490 & -30.7\% \\
 & Q4b SLD/EM, 보정셋 제외 학습 & none & 0.9870 & +0.3728 & [0.950--1.000] & 0.3701 & -39.2\% \\
 & Q5 SLD/EM + Platt 보정 & training labels only & 1.0000 & +0.3858 & [0.811--1.000] & 0.3858 & -45.4\% \\
 & Q6 ACC/BBSE-hard & training labels only & 0.8945 & +0.2803 & [0.000--1.000] & 0.3179 & -18.2\% \\
 & Q7 BBSE-soft & training labels only & 0.5257 & -0.0886 & [0.142--0.907] & 0.2210 & 20.7\% \\
 & Q8 HDy & training labels only & 1.0000 & +0.3858 & [0.970--1.000] & 0.3858 & -45.4\% \\
 & Q9 MS & training labels only & 0.4620 & -0.1522 & [0.019--1.000] & 0.2064 & 26.6\% \\
 & Q0 참값 & validation labels & 0.6142 & +0.0000 & [0.558--0.670] & 0.2438 & 11.5\% \\
\botrule
\end{tabular}
\footnotetext{Sensitivities, specificities and score distributions are estimated on a held-out slice of the training era, so no estimator uses validation labels except the oracle row. Intervals are 2.5--97.5 percentiles of a bootstrap of the validation set with the predictions held fixed. Recovery is defined as elsewhere in this work, relative to an oracle Platt upper bound. \textbf{For DILIrank every estimator overshoots, every recovery is negative, and no interval covers the truth.} The published remedy of calibrating the classifier before running expectation maximisation does not help in either dataset. For DICTrank, 3 estimators do give a positive recovery, but all of them undershoot the prevalence, and undershooting is already known to lower ECE below the oracle value, so the gain is plausibly a by-product of the error rather than a successful estimate; their intervals are too wide to identify the prevalence at all. Probabilistic adjusted count and soft black-box shift estimation are algebraically the same estimator for a binary outcome, so their agreement is a check rather than independent evidence.}
\end{table}

\begin{table}[h]
\caption{Removing the arbitrariness from the point estimates of the recalibration and conformal analyses}\label{tab:tieavg}
\begin{tabular}{@{}llcccccc@{}}
\toprule
Dataset & Quantity & Mean & SD & 2.5--97.5 pct & Min--max & Reported & Distance \\
\midrule
DILIrank & Uncalibrated ECE & 0.4419 & 0.0076 & 0.4295--0.4557 & 0.4293--0.4577 & 0.4466 & +0.6\,SD \\
 & Uncalibrated ECE, calibration set held out & 0.4596 & 0.0051 & 0.4510--0.4690 & 0.4504--0.4695 & 0.4577 & -0.4\,SD \\
 & Oracle Platt ECE (upper bound) & 0.0446 & 0.0162 & 0.0198--0.0869 & 0.0160--0.0912 & 0.0482 & +0.2\,SD \\
 & Platt ECE & 0.4223 & 0.0048 & 0.4130--0.4312 & 0.4112--0.4362 & 0.4147 & -1.6\,SD \\
 & Isotonic ECE & 0.4280 & 0.0066 & 0.4156--0.4388 & 0.4125--0.4417 & 0.4085 & -3.0\,SD \\
 & Prior-shift ECE & 0.2966 & 0.0099 & 0.2824--0.3156 & 0.2766--0.3285 & 0.2959 & -0.1\,SD \\
 & Platt recovery (\%) & 4.89 & 2.38 & 0.66--8.77 & 0.14--9.63 & 8.01 & +1.3\,SD \\
 & Isotonic recovery (\%) & 3.46 & 2.44 & -0.91--8.16 & -1.10--10.83 & 9.58 & +2.5\,SD \\
 & Prior-shift recovery (\%) & 36.64 & 3.06 & 30.63--42.37 & 29.65--44.39 & 37.82 & +0.4\,SD \\
 & Conformal 90\% coverage & 0.6699 & 0.0305 & 0.6168--0.7229 & 0.5897--0.7436 & 0.7372 & +2.2\,SD \\
 & Conformal 80\% coverage & 0.4897 & 0.0133 & 0.4679--0.5128 & 0.4679--0.5192 & 0.4808 & -0.7\,SD \\
 & AUROC & 0.6463 & 0.0085 & 0.6325--0.6637 & 0.6300--0.6696 & 0.6603 & +1.6\,SD \\
\midrule
DICTrank & Uncalibrated ECE & 0.2720 & 0.0064 & 0.2608--0.2829 & 0.2584--0.2882 & 0.2726 & +0.1\,SD \\
 & Uncalibrated ECE, calibration set held out & 0.2945 & 0.0070 & 0.2809--0.3053 & 0.2796--0.3160 & 0.2942 & -0.0\,SD \\
 & Oracle Platt ECE (upper bound) & 0.0323 & 0.0116 & 0.0154--0.0537 & 0.0154--0.0591 & 0.0234 & -0.8\,SD \\
 & Platt ECE & 0.1741 & 0.0046 & 0.1650--0.1829 & 0.1627--0.1833 & 0.1679 & -1.3\,SD \\
 & Isotonic ECE & 0.1829 & 0.0063 & 0.1688--0.1940 & 0.1662--0.1993 & 0.1739 & -1.4\,SD \\
 & Prior-shift ECE & 0.2270 & 0.0082 & 0.2131--0.2424 & 0.2085--0.2438 & 0.2438 & +2.1\,SD \\
 & Platt recovery (\%) & 40.89 & 3.18 & 36.31--48.24 & 35.18--49.28 & 41.99 & +0.3\,SD \\
 & Isotonic recovery (\%) & 37.20 & 3.51 & 32.78--44.42 & 29.14--46.04 & 39.60 & +0.7\,SD \\
 & Prior-shift recovery (\%) & 18.83 & 3.66 & 12.81--25.77 & 12.32--27.37 & 11.55 & -2.0\,SD \\
 & Conformal 90\% coverage & 0.8168 & 0.0158 & 0.7903--0.8456 & 0.7790--0.8502 & 0.8240 & +0.5\,SD \\
 & Conformal 80\% coverage & 0.6942 & 0.0152 & 0.6667--0.7257 & 0.6629--0.7303 & 0.6966 & +0.2\,SD \\
 & AUROC & 0.6122 & 0.0083 & 0.5969--0.6262 & 0.5929--0.6265 & 0.6045 & -0.9\,SD \\
\botrule
\end{tabular}
\footnotetext{Three arbitrary implementation choices are randomised jointly over 50 repetitions: the assignment of boundary-year ties to the internal calibration set, the row order of the fitting set, and the row order of the full-training model. The last matters because the denominator of the recovery rate comes from the full-training model, so randomising only the numerator would measure half of the arbitrariness. The mean does not depend on any of those choices and is therefore the estimate to read; \emph{Reported} is the single realisation printed in the main tables, and \emph{Distance} expresses its offset from the mean in repetition standard deviations. 22 of 24 reported values fall inside the min--max range of the 50 repetitions. Every conclusion holds in all repetitions: nominal coverage is missed in 50/50, within-era recalibration recovers less than the oracle in 50/50, and the best-performing method differs between the two datasets in 50/50.}
\end{table}

\begin{table}[h]
\caption{A per-class conformal guarantee recovers coverage in one dataset only}\label{tab:mondrian}
\begin{tabular}{@{}llllccccc@{}}
\toprule
Dataset & Target & Split & Procedure & Coverage & Positive & Negative & Set size & $z_{+}$ & $z_{-}$ \\
\midrule
DILIrank & 90\% & Temporal & Marginal & 0.7372 & 0.9583 & 0.6389 & 1.49 & +1.35 & -9.05 \\
DILIrank & 90\% & Temporal & Mondrian & 0.8654 & 0.9167 & 0.8426 & 1.60 & +0.38 & -1.99 \\
DILIrank & 90\% & Random & Marginal & 0.9087 $\pm$ 0.0279 & 0.9531 & 0.8450 & 1.59 & --- & --- \\
DILIrank & 90\% & Random & Mondrian & 0.9186 $\pm$ 0.0333 & 0.9136 & 0.9255 & 1.65 & --- & --- \\
DILIrank & 90\% & Random, prevalence matched & Marginal & 0.8786 $\pm$ 0.0417 & 0.9542 & 0.8450 & 1.61 & --- & --- \\
DILIrank & 90\% & Random, prevalence matched & Mondrian & 0.9202 $\pm$ 0.0398 & 0.9080 & 0.9255 & 1.63 & --- & --- \\
DILIrank & 80\% & Temporal & Marginal & 0.4808 & 0.9167 & 0.2870 & 1.04 & +2.02 & -13.33 \\
DILIrank & 80\% & Temporal & Mondrian & 0.7628 & 0.8333 & 0.7315 & 1.38 & +0.58 & -1.78 \\
DILIrank & 80\% & Random & Marginal & 0.8151 $\pm$ 0.0336 & 0.8907 & 0.7091 & 1.32 & --- & --- \\
DILIrank & 80\% & Random & Mondrian & 0.8266 $\pm$ 0.0310 & 0.8233 & 0.8314 & 1.38 & --- & --- \\
DILIrank & 80\% & Random, prevalence matched & Marginal & 0.7654 $\pm$ 0.0512 & 0.8922 & 0.7091 & 1.34 & --- & --- \\
DILIrank & 80\% & Random, prevalence matched & Mondrian & 0.8267 $\pm$ 0.0368 & 0.8162 & 0.8314 & 1.36 & --- & --- \\
\midrule
DICTrank & 90\% & Temporal & Marginal & 0.8240 & 0.9939 & 0.5534 & 1.43 & +4.01 & -11.73 \\
DICTrank & 90\% & Temporal & Mondrian & 0.9213 & 0.9573 & 0.8641 & 1.76 & +2.45 & -1.22 \\
DICTrank & 90\% & Random & Marginal & 0.9062 $\pm$ 0.0277 & 0.9827 & 0.7020 & 1.48 & --- & --- \\
DICTrank & 90\% & Random & Mondrian & 0.9195 $\pm$ 0.0318 & 0.9216 & 0.9128 & 1.70 & --- & --- \\
DICTrank & 90\% & Random, prevalence matched & Marginal & 0.8734 $\pm$ 0.0390 & 0.9809 & 0.7020 & 1.50 & --- & --- \\
DICTrank & 90\% & Random, prevalence matched & Mondrian & 0.9201 $\pm$ 0.0361 & 0.9247 & 0.9128 & 1.70 & --- & --- \\
DICTrank & 80\% & Temporal & Marginal & 0.6966 & 0.9878 & 0.2330 & 1.15 & +6.01 & -14.39 \\
DICTrank & 80\% & Temporal & Mondrian & 0.8464 & 0.8720 & 0.8058 & 1.55 & +2.30 & +0.15 \\
DICTrank & 80\% & Random & Marginal & 0.8182 $\pm$ 0.0337 & 0.9466 & 0.4787 & 1.21 & --- & --- \\
DICTrank & 80\% & Random & Mondrian & 0.8219 $\pm$ 0.0313 & 0.8231 & 0.8177 & 1.42 & --- & --- \\
DICTrank & 80\% & Random, prevalence matched & Marginal & 0.7657 $\pm$ 0.0407 & 0.9458 & 0.4787 & 1.22 & --- & --- \\
DICTrank & 80\% & Random, prevalence matched & Mondrian & 0.8265 $\pm$ 0.0359 & 0.8321 & 0.8177 & 1.42 & --- & --- \\
\botrule
\end{tabular}
\footnotetext{Marginal split conformal calibrates one quantile on the whole calibration set and guarantees only overall coverage, so the per-class figures of the main conformal analysis are descriptive. Mondrian conformal calibrates a separate quantile on each true-label stratum of the calibration set, which guarantees coverage within each class and needs no knowledge of the test label. Because it relies only on exchangeability within a class, it stays valid under arbitrary label shift, which makes it a diagnostic as well as a remedy: whatever coverage deficit survives it cannot be attributed to the change in prevalence. The random split is a size-matched control over 20 repetitions, and the prevalence-matched control subsamples each repetition's evaluation set to the positive rate of the temporal validation set, which changes the label marginal and nothing else. Two checks confirm the implementation: under the exchangeable control Mondrian attains nominal coverage in both classes, and moving to the prevalence-matched control moves overall coverage by 0.0301--0.0525 for the marginal procedure but by at most 0.0046 for Mondrian. $z$ is a one-sided binomial statistic against the nominal level and is reported only for the single temporal split. Under the temporal split Mondrian attains nominal coverage in both classes for DICTrank, while DILIrank still falls significantly short in the negative class, so that dataset carries a coverage deficit that label shift does not explain. The cost is set size, which rises by 0.11--0.40 under the temporal split; no prediction set is empty.}
\end{table}

\begin{table}[h]
\caption{Every test previously outside a correction family, folded back in}\label{tab:fdrall}
\begin{tabular}{@{}lcccc@{}}
\toprule
Analysis & Tests & Significant in F4 & Smallest $q$ (F4) & Smallest $q$ (F0) \\
\midrule
Label-conditional conformal coverage & 16 & 4 & $<$0.0001 & $<$0.0001 \\
Reversed temporal split & 6 & 4 & 0.0315 & 0.0256 \\
Non-temporal axis splits & 10 & 2 & 0.0315 & 0.0256 \\
Concept shift & 4 & 0 & 0.0537 & 0.0531 \\
Applicability domain & 9 & 0 & 0.0575 & 0.0606 \\
Label binarisation & 12 & 8 & 0.0315 & 0.0306 \\
\midrule
\textbf{All newly included} & \textbf{57} & \textbf{18} & $<$0.0001 & $<$0.0001 \\
Pre-existing families & 20 & 8 (unchanged families) & --- & 0.0256 \\
\botrule
\end{tabular}
\footnotetext{The main multiple-comparison analysis corrects within three pre-specified families and leaves the tests of several later sections outside any of them, on the stated grounds that no confirmatory claim rests on those p-values. That is a weak position to defend, so the 57 excluded tests are folded back in here: F4 applies Benjamini--Hochberg to those tests as a single family, and F0 applies it to all 77 tests in the paper at once, which is the most conservative construction available. Negative-control rows are kept rather than dropped, since a larger family makes the correction stricter. Coverage $z$ statistics are converted to one-sided p-values under a normal approximation. 18 of the 57 newly included tests survive F4. Under F0, 1 of the 8 previously significant results loses significance (DICTrank, R3. 무작위 + 학습·검증 비율맞춤), and that result was already flagged as resting on its effect size rather than its p-value. The one claim that does not survive correction is the residual coverage deficit of label-conditional conformal prediction, which is therefore reported as an effect size throughout.}
\end{table}

\begin{table}[h]
\caption{Calibration loss is present in every subgroup that could be measured}\label{tab:subcalib}
\begin{tabular}{@{}llccccccc@{}}
\toprule
Dataset & Stratum & $n$ & Positive rate & AUROC & ECE & Control ECE & ECE ratio & 95\% CI \\
\midrule
DILIrank & ATC classes other than L & 118 & 0.1780 & 0.6490 & 0.5405 & 0.3572 & \textbf{1.51} & [1.315--1.751] \\
 & Molecular weight $>$ 500 & 45 & 0.3556 & 0.6789 & 0.3775 & 0.2634 & \textbf{1.43} & [1.005--2.016] \\
 & Molecular weight $\leq$ 500 & 111 & 0.2883 & 0.6424 & 0.4822 & 0.2818 & \textbf{1.71} & [1.432--2.049] \\
 & Fewer than four aromatic rings & 124 & 0.2661 & 0.6593 & 0.4531 & 0.2978 & \textbf{1.52} & [1.292--1.826] \\
 & At least six nitrogen atoms & 40 & 0.4000 & 0.7812 & 0.3658 & 0.2407 & \textbf{1.52} & [1.016--2.135] \\
 & Fewer than six nitrogen atoms & 116 & 0.2759 & 0.5893 & 0.4765 & 0.2915 & \textbf{1.63} & [1.390--1.972] \\
 & logP $\geq$ 3 & 77 & 0.3766 & 0.5553 & 0.4717 & 0.2267 & \textbf{2.08} & [1.666--2.641] \\
 & logP $<$ 3 & 79 & 0.2405 & 0.7132 & 0.4343 & 0.3030 & \textbf{1.43} & [1.149--1.769] \\
\midrule
DICTrank & ATC class L (antineoplastic and immunomodulating) & 94 & 0.7553 & 0.5903 & 0.1676 & 0.1333 & 1.26 & [0.774--1.942] \\
 & ATC classes other than L & 167 & 0.5449 & 0.6385 & 0.3345 & 0.2440 & \textbf{1.37} & [1.121--1.696] \\
 & Molecular weight $>$ 500 & 98 & 0.6224 & 0.5162 & 0.2752 & 0.1983 & \textbf{1.39} & [1.060--1.967] \\
 & Molecular weight $\leq$ 500 & 169 & 0.6095 & 0.6840 & 0.2892 & 0.1958 & \textbf{1.48} & [1.165--1.841] \\
 & At least four aromatic rings & 63 & 0.6825 & 0.5860 & 0.2830 & 0.1796 & \textbf{1.58} & [1.095--2.191] \\
 & Fewer than four aromatic rings & 204 & 0.5931 & 0.6300 & 0.2930 & 0.2173 & \textbf{1.35} & [1.101--1.683] \\
 & At least six nitrogen atoms & 84 & 0.7024 & 0.5166 & 0.2387 & 0.1599 & \textbf{1.49} & [1.088--2.200] \\
 & Fewer than six nitrogen atoms & 183 & 0.5738 & 0.6740 & 0.3035 & 0.2292 & \textbf{1.32} & [1.090--1.649] \\
 & logP $\geq$ 3 & 152 & 0.6645 & 0.5791 & 0.2371 & 0.1617 & \textbf{1.47} & [1.187--2.029] \\
 & logP $<$ 3 & 115 & 0.5478 & 0.6514 & 0.3273 & 0.2388 & \textbf{1.37} & [1.065--1.802] \\
\botrule
\end{tabular}
\footnotetext{The subgroup analyses in the main text report discrimination only, which leaves open whether the calibration loss is concentrated in a therapeutic or structural class; if it were, the practical remedy would be to withhold probabilities for that class. Each stratum is therefore paired with random splits matched on both stratum size and positive rate, over 30 repetitions, following the design used for the applicability-domain analysis. Intervals are percentiles of a within-stratum bootstrap and do not include the variation of the control mean, so they are narrower than the true intervals. The underpowered rule used elsewhere in this work (fewer than 40 compounds, or fewer than ten in the minority class) removes 27 of 46 strata, which excludes almost every individual ATC class: the largest class in the DILIrank validation set holds 32 compounds. Of the 18 strata that remain, 17 show a ratio above one whose interval excludes one, with ratios from 1.32 to 2.08. The loss is therefore not concentrated anywhere, and no subgroup-based rule for withholding probabilities follows from it.}
\end{table}

\begin{table}[h]
\caption{Retraining at all matters more than how often}\label{tab:interval}
\begin{tabular}{@{}lllcccccc@{}}
\toprule
Dataset & Recalibrated & Interval & Refits & Mean staleness & ECE & AUROC & $\Delta$ECE vs annual & 95\% CI \\
\midrule
DILIrank & yes & 1 years & 13 & 0.00 & 0.2836 & 0.7218 & reference & --- \\
 & yes & 2 years & 7 & 0.53 & 0.3029 & 0.7141 & \textbf{+0.0186} & [0.005--0.027] \\
 & yes & 3 years & 5 & 0.99 & 0.3054 & 0.7307 & \textbf{+0.0232} & [0.010--0.035] \\
 & yes & 5 years & 3 & 1.89 & 0.3292 & 0.7247 & \textbf{+0.0449} & [0.031--0.056] \\
 & yes & never (fixed) & 1 & 5.71 & 0.4224 & 0.6632 & \textbf{+0.1387} & [0.115--0.162] \\
 & no & 1 years & 13 & 0.00 & 0.3751 & 0.6842 & reference & --- \\
 & no & 2 years & 7 & 0.53 & 0.3662 & 0.6823 & -0.0041 & [-0.024, 0.016] \\
 & no & 3 years & 5 & 0.99 & 0.3779 & 0.6736 & +0.0053 & [-0.020, 0.027] \\
 & no & 5 years & 3 & 1.89 & 0.3858 & 0.6867 & +0.0126 & [-0.017, 0.039] \\
 & no & never (fixed) & 1 & 5.71 & 0.4384 & 0.6476 & \textbf{+0.0681} & [0.033--0.105] \\
\midrule
DICTrank & yes & 1 years & 14 & 0.00 & 0.0960 & 0.5943 & reference & --- \\
 & yes & 2 years & 7 & 0.55 & 0.1032 & 0.6064 & +0.0038 & [-0.016, 0.020] \\
 & yes & 3 years & 6 & 1.16 & 0.1062 & 0.5809 & +0.0066 & [-0.023, 0.034] \\
 & yes & 5 years & 4 & 1.86 & 0.1199 & 0.5948 & +0.0185 & [-0.008, 0.039] \\
 & yes & never (fixed) & 1 & 6.43 & 0.1716 & 0.6259 & \textbf{+0.0726} & [0.039--0.096] \\
 & no & 1 years & 14 & 0.00 & 0.2583 & 0.6227 & reference & --- \\
 & no & 2 years & 7 & 0.55 & 0.2574 & 0.6218 & -0.0006 & [-0.019, 0.019] \\
 & no & 3 years & 6 & 1.16 & 0.2532 & 0.6284 & +0.0008 & [-0.023, 0.026] \\
 & no & 5 years & 4 & 1.86 & 0.2538 & 0.6241 & -0.0016 & [-0.027, 0.022] \\
 & no & never (fixed) & 1 & 6.43 & 0.2702 & 0.6198 & +0.0200 & [-0.008, 0.049] \\
\botrule
\end{tabular}
\footnotetext{Each row predicts every validation compound with the model from the most recent refit at or before its approval year, so the evaluation set is the same in every row and the differences are paired on the same compounds. Mean staleness is the average number of years between a compound's approval and the close of the training window that produced its prediction; a five-year interval gives a mean staleness of about two years, not five, because compounds are spread across the cycle. Recalibration repeats the Platt step of the static analysis at each refit, on the most recent quarter of the training window. Differences are paired bootstrap estimates against annual retraining. With recalibration, never refitting is worse than refitting annually by +0.1387 and +0.0726, both significant, while stretching the interval to five years costs only +0.0449 and +0.0185 -- 3.1 and 3.9 times less. The staleness gradient is significant in DILIrank but in DICTrank no interval up to five years is distinguishable from annual. Without recalibration no gradient is detectable at all, which is the same conclusion the single-interval analysis reached: the recalibration step is what does the work. These rows assume a label is available in the year of approval and therefore bound rather than estimate the benefit.}
\end{table}
